\documentclass[11pt]{article}
\usepackage[T1]{fontenc}
\usepackage{amsmath,amssymb}
\usepackage{geometry}
\geometry{margin=1in}
\title{Zoom-in glass initial conditions on Poincar\'e dodecahedral space}
\author{StePS}
\date{}
\begin{document}
\maketitle

\section{Scope and conventions}
The generator constructs a \emph{pre-glass} particle load on the quotient
\(\mathcal M=S^3/I^*\), with a refined spherical region, a transition, a
constant-resolution buffer, and a homogeneous exterior.  The subsequent
reverse-gravity relaxation and cosmological perturbations are separate
calculations.  Distances are comoving Mpc, particle masses are in
\(10^{11}M_\odot\), and the scale factor used for the volume and density is
\(a=1\).

Let \(R\) be the curvature radius of the covering \(S^3\), and \(q\) a unit
quaternion representing a point. StePS uses the stereographic coordinate
\(\boldsymbol{x}=R\boldsymbol{q}/(1+q_0)\), with inverse
\[
 q_0=\frac{R^2-|\boldsymbol{x}|^2}{R^2+|\boldsymbol{x}|^2},
 \qquad
 \boldsymbol q=\frac{2R\boldsymbol{x}}{R^2+|\boldsymbol{x}|^2}.
\]
The fundamental domain is the Voronoi cell of \(e=(1,0,0,0)\), namely
\(q\cdot e\ge q\cdot g\) for every \(g\in I^*\).  Its volume is
\(V_{\rm PDS}=2\pi^2R^3/120=\pi^2R^3/60\).  The code samples the existing
homogeneous load in this domain, rather than a Euclidean dodecahedron.

\section{Zoom centre and injectivity}
The centre \(q_c\) may be the domain centre or a point specified by three
stereographic coordinates inside the cell.  The quotient distance from \(q\)
to \(q_c\) is
\[
 \chi(q,q_c)=\arccos\!\left[\max_{g\in I^*} q\cdot(gq_c)\right].
\]
The group acts by left quaternion multiplication.  Its transformations are
Clifford translations, and the shortest non-identity displacement has
angular length \(\pi/5\). Thus the injectivity radius at every centre is
\(R\pi/10\). The code computes this from the group and requires
\(\chi_{\rm zoom}<\pi/10\). By default it additionally requires
\(\chi_{\rm zoom}\le0.8\,\pi/10\). The expert override only relaxes this
safety margin; it never permits the ball to reach the cut locus.

An off-centre ball can cross the chosen \emph{coordinate} cell's faces without
crossing the quotient cut locus. Its generated points are reduced to the
fundamental domain using the same nearest-group-element and left-multiplication
convention as StePS. All inside/outside tests use the quotient distance, not a
Euclidean distance in the stereographic chart.

\section{Radial geometry and the map}
With geodesic polar angle \(\chi\) around \(q_c\),
\[
 ds^2=R^2(d\chi^2+\sin^2\chi\,d\Omega_2^2),\qquad
 V(<\chi)=2\pi R^3 F(\chi),\qquad
 F(\chi)=\chi-\tfrac12\sin 2\chi .
\]
The auxiliary radius is \(\omega\), and the central M\"obius map is
\[
 f_M(\omega)=2\arctan\!\left[\Lambda\tan(\omega/2)\right],\qquad
 f'_M(\omega)=\Lambda\frac{1+\tan^2(\omega/2)}
 {1+\Lambda^2\tan^2(\omega/2)},\quad 0<\Lambda<1.
\]
Near the centre its linear scale is \(\Lambda\) and its volume Jacobian is
\(\Lambda^3\). A nontrivial map that is the identity outside a finite ball
cannot simultaneously have a nondecreasing Jacobian below unity inside that
ball: its integrated volume would be smaller than the original ball's volume.
The former identity-matched construction therefore produced a Jacobian above
unity followed by a decline. Instead, we match the \emph{resolution}, not
the auxiliary coordinate, to the unchanged physical exterior.

Let \(\chi_1\) and \(\chi_2\) be the physical angular radii of the core and
transition, given by the configured fractions of \(\chi_{\rm zoom}\). In
physical coordinates the pure M\"obius Jacobian is
\[
 J_M(\chi)=\left[\Lambda\cos^2(\chi/2)
                  +\Lambda^{-1}\sin^2(\chi/2)\right]^3.
\]
For \(u=(\chi-\chi_1)/(\chi_2-\chi_1)\), clipped to \([0,1]\), define
\(S(u)=6u^5-15u^4+10u^3\) and
\[
 J_{\rm target}(\chi)=[1-S(u)]J_M(\chi)+S(u).
\]
We require \(J_M(\chi_2)<1\). Then \(J_{\rm target}\) is nondecreasing,
equals the M\"obius Jacobian in the core, and is one in the buffer. The map
\(\chi=f(\omega)\) is obtained from
\[
 \frac{df}{d\omega}=J_{\rm target}(f)
 \left[\frac{\sin\omega}{\sin f}\right]^2,
 \qquad f(\omega_1)=\chi_1,\qquad
 \omega_1=2\arctan\!\left[\frac{\tan(\chi_1/2)}{\Lambda}\right].
\]
The core uses \(f_M\) exactly. The remainder is integrated numerically and
interpolated with a cubic Hermite spline. The code checks positivity of the
interpolated derivative, including interior extrema. In the buffer \(J=1\),
but the auxiliary radius is not the physical radius: their volume coordinates
differ by a constant. In particular \(\omega_{\rm zoom}>\chi_{\rm zoom}\).
Exterior positions are independently sampled in the physical PDS domain.

The physical-to-auxiliary volume Jacobian is
\[
 J(\omega)=f'(\omega)
 \left[\frac{\sin f(\omega)}{\sin\omega}\right]^2,\qquad
 J(0)=\Lambda^3 .
\]
For equal auxiliary-volume shells with the same particle count, the shell
mass is proportional to the shell-averaged \(J\). For equal-\(\omega\)
shells the auxiliary volumes themselves vary with radius, producing an
additional mass hierarchy; the central mass need not be \(\Lambda^3 m_{\rm out}\).

\section{Shell construction}
Shells are equally spaced either in \(\omega\) or in the auxiliary volume
coordinate \(F(\omega)\). Their physical boundaries are
\(\chi_i=f(\omega_i)\). The exact physical shell volume is
\(V_i=2\pi R^3[F(\chi_{i+1})-F(\chi_i)]\). The radius \(\chi_{c,i}\)
reported in diagnostics is the equal-volume midpoint,
\[
 F(\chi_{c,i})=\tfrac12[F(\chi_i)+F(\chi_{i+1})].
\]
This is the median in shell volume, not the mean radial coordinate.
Each shell has \(N_{\rm shell}\) independent uniformly random directions, like the
\(R^3\) glass generator. Radii are randomly sampled with the \(S^3\) volume
measure, rather than uniformly in radius; thus the expected mass density
within each shell is uniform. For each direction \(\hat n\) and radius \(\chi\),
the point on \(S^3\) is
\[
 q=q_c\bigl(\cos\chi,\ \sin\chi\hat n\bigr).
\]
Each retained shell contains exactly \(N_{\rm shell}\) particles, with
\(N_{\rm bins}\) specifying the original radial binning. For all but the last
shell, their mass is
\(m_i=\bar\rho V_i/N_{\rm shell}\). The last shell shares a mass with the
exterior, as described below. The code checks that shell masses do not
decrease outwards and, if configured, enforces an adjacent mass-ratio limit.
At least two whole shells must lie in the buffer. With equal-\(\omega\)
binning the second shell can be approximately seven times as massive as the
first; an optional ratio limit must allow for this.

\section{Constant-resolution central region}
An optional central ball of geodesic radius \(r_c=R\chi_c\) replaces the
innermost shells. Its particles are independently sampled uniformly in the
\(S^3\) volume, wrapped to the PDS fundamental domain, and all have the
same mass. If \(V_c=V(<\chi_c)\) and \(m_1\) is the mass of the first
retained shell, choose
\[
 N_c=\left\lceil\frac{\bar\rho V_c}{m_1}\right\rceil,\qquad
 m_c=\frac{\bar\rho V_c}{N_c}.
\]
Thus the central mass and density are exact, and \(m_c\le m_1\);
integer counts can leave a mass step at the join.

The radius can be specified directly in Mpc. It must be smaller than the zoom
radius; the remaining shells are re-spaced from the exact selected radius to
the zoom boundary, preserving their configured binning and \(N_{\rm shell}\).
Alternatively, \texttt{auto} scans the original shells from outside inward and
selects the boundary outside the outermost pair for which the outer-to-inner
particle-mass ratio exceeds the configured threshold (default \(1.25\)).
If integer rounding makes the central-to-shell ratio too large, the boundary
is moved outwards until the threshold is met. If no pair exceeds the threshold,
no central ball is added. The automatic mode rejects cases where its bound
cannot be met while retaining the required buffer shells.

\section{Exterior and mass conservation}
Let \(V_{\rm ext}=V_{\rm PDS}-V(<\chi_{\rm zoom})\) and \(V_{\rm last}\)
be the final shell volume. We select the integer exterior particle count
\[
 N_{\rm ext}=\left\lfloor N_{\rm shell}
 \frac{V_{\rm ext}}{V_{\rm last}}\right\rfloor .
\]
It is not generally possible for an integer \(N_{\rm ext}\) to give both the
last shell and exterior exactly the same particle mass and exactly
\(\bar\rho\) in each region. We assign both regions the common mass
\[
 m_{\rm out}=\bar\rho\,
 \frac{V_{\rm PDS}-V_c-\sum_{i<{\rm last}}V_i}
 {N_{\rm ext}+N_{\rm shell}}.
\]
This preserves \(\sum m=\bar\rho V_{\rm PDS}\) and removes the jump at the
join. Flooring ensures that the correction does not lower the edge mass
below its nominal shell value. The final shell and exterior densities can
differ slightly from \(\bar\rho\); the relative correction is recorded and
reported, with a warning above one per cent. The exterior points are sampled
uniformly from the PDS domain outside the quotient zoom ball. Finite-sampling
fluctuations remain locally. The total particle count is
\(N_c+N_{\rm retained}N_{\rm shell}+N_{\rm ext}\); \texttt{NPART} is ignored in zoom
mode and remains the particle count in homogeneous mode.

\section{Configuration and diagnostics}
The example \texttt{GlassIC\_PDS\_zoom.yaml} enables the
\texttt{zoom\_glass} mapping. The zoom radius can be given as a fraction of
the injectivity radius or in Mpc, but not both. The configuration also sets
\(\Lambda\), physical core and transition fractions, shell count and binning,
random seed, \texttt{NRBINS}, \texttt{NSHELL}, buffer count, and an optional
maximum adjacent-shell mass ratio, and the optional central radius and automatic
ratio threshold.
The sampler is \texttt{random}; the former \texttt{fibonacci} setting is
accepted with a warning and also uses random sampling. The old shape and mass weights, minimum angular count, and shell-rotation
setting no longer affect sampling or counts.
The output extension selects StePS HDF5 or ASCII format. The original
homogeneous mode is unchanged when zoom is disabled.

Zoom particles are saved in increasing azimuth \(\phi=\operatorname{atan2}(y,x)\)
in the zoom centre's local quaternion frame. For each particle the nearest
PDS image to the zoom centre supplies the local direction, so the definition
also applies across identified faces. HDF5 particle IDs increase with this
order; ASCII rows use the same order. Equal-distance image ties are resolved by
the group ordering. Sorting requires additional memory beyond positions and
masses, as reported by the generator and resolution calculator.

Each zoom output has three sidecars with the output filename's extension
removed: \texttt{.zoom.json} for parameters, counts and mass accounting;
\texttt{.zoom-shells.csv} for shell boundaries, masses, counts, spacings,
and Jacobian; and \texttt{.zoom-map.csv} for the map,
derivative and Jacobian. Angles in these tables are dimensionless radians;
lengths are in Mpc. The random exterior and shell samples have separately
recorded seeds.

\section{Limitations and validation}
The shell construction deliberately introduces a radial pattern that must be
removed by glass relaxation. Exact shell masses do not guarantee absence of
noise, a smooth nearest-neighbour distribution, or a stable post-relaxation
zoom profile. The configured mass-ratio check covers adjacent shells and the
exterior join, not every spatially neighbouring pair. Check pre- and
post-relaxation density, particle spacing, mass segregation, shell structure
and boundary crossings before scientific use. For the zoom load the generator
estimates the minimum-mass particle separation as
\(d_{\min}=(m_{\min}/\bar\rho)^{1/3}\) and recommends \(d_{\min}/40\) as a
softening length. This density-based estimate is not a measurement of actual
nearest-neighbour distances, nor a guarantee of optimal softening. The code
does not perform or certify long-term relaxation.
\end{document}
